%! Representing a Quantitative Variable with Graphs — Unit 1, Lesson 1.3 — Guided Notes & Practice (Answer Key)
% MIRROR of ../notes/main.tex: -key in place of -boxes, \termblank -> \termans, \blank -> \ans, and the
% answer argument of every \pcell, \writespace and \labelbox filled. Nothing else differs.
% THE in-class document, in the MAIN IDEAS / NOTES shape (LESSON_SHAPE.md §1): the vocabulary
% box, then ONE two-column guidednotes table. Each row is one idea — a label on the left; on the
% right a SHORT printed definition, ONE large pre-drawn display, then a grid of a few numbered
% problems with real writing room. Problems are numbered 1..N through the component.
% DENSITY: a \blank{} only where a word or number IS the answer (the interval table, naming a
% display); no mid-sentence blanks; two problems across; 2–3cm of answer space each.
%   I DO   : the teacher reads the definition, marks up the display, works the first problem.
%   WE DO  : the class works the rest of each grid, students holding the pen; the last row,
%            Guided Practice, is worked entirely together in a fresh context.
%   YOU DO : nothing on this page — ../homework/ IS the individual practice.
% The crux is the DRAWN TWICE row — the same 24 times at two widths tell two different stories,
% and neither picture is wrong (EK UNC-1.G.1).
%
% THE DATA (24 quiz finishing times):
%   11 12 12 14 15 15 16 17 18 18 19 20 | 34 35 36 36 38 39 40 41 42 43 44 46
% Width-5 bins from 10: 4, 7, 1, 0, 1, 5, 5, 1.   Width-20 bins from 10: 12, 12.
%
% PAGE LOCKSTEP with notes_key/: \termblank -> \termans, \blank -> \ans, and the answer argument
% of every \pcell, \writespace and \labelbox filled. Nothing else differs. A row cannot break
% across pages: figure and grid are separate sub-rows (`\\` then `& ...`).
\documentclass[12pt]{article}
\usepackage{apstats-article}
\usepackage{apstats-key}   % pulls in -boxes; do NOT also load -boxes

\begin{document}

\pageheader{Unit 1, Lesson 1.3}{Guided Notes \& Practice --- Answer Key}
% NAMESTRIP: no \namedateperiod here — the name row lives on the cover only.

\vspace{2pt}

\boxguard[16]
\begin{vocabbox}
\small
Fill in each term \textbf{as we name it} in the notes below.
\par
\termans{Discrete variable}{counted; only certain values can happen, nothing sits between neighbours.}
\par
\termans{Continuous variable}{measured; any value in an interval can happen.}
\par
\termans{Histogram}{grouped quantitative data: one bar per interval, height = count, bars touching.}
\par
\termans{Interval width}{the size of each histogram interval; changing it changes the shape, not the data.}
\par
\end{vocabbox}

\small
\begin{guidednotes}
% ── ROW 1: discrete vs continuous ────────────────────────────────────────────
\mainidea[Discrete \&]{Continuous} &
A \textbf{discrete} variable is \emph{counted}: only certain values can happen, and nothing sits
between two neighbours. A \textbf{continuous} variable is \emph{measured}: any value in an
interval can happen. \textbf{The test:} could a finer instrument split the value?

\vspace{4pt}
\begin{center}
\begin{tikzpicture}[scale=0.95]
  \node[left, font=\small\bfseries\color{navy}] at (-0.4,1.3) {Discrete};
  \draw[->, thick] (-0.2,1.3) -- (7.2,1.3);
  \foreach \x in {0,...,6} {
    \draw (\x,1.22) -- (\x,1.38); \node[below, font=\tiny] at (\x,1.15) {\x};
    \fill[navy] (\x,1.3) circle (0.09);}
  \node[font=\tiny] at (3.5,0.62) {\emph{siblings} --- only these values can happen};
  \node[left, font=\small\bfseries\color{navy}] at (-0.4,-0.6) {Continuous};
  \fill[navy!25] (0,-0.72) rectangle (6.6,-0.48);
  \draw[->, thick] (-0.2,-0.6) -- (7.2,-0.6);
  \foreach \x/\l in {0/0, 2/4, 4/8, 6/12} {
    \draw (\x,-0.68) -- (\x,-0.52); \node[below, font=\tiny] at (\x,-0.75) {\l};}
  \draw[goldacc, very thick] (1.75,-0.95) -- (1.75,-0.25);
  \node[above, font=\tiny] at (1.75,-0.25) {3.5 min};
  \node[font=\tiny] at (3.5,-1.4) {\emph{minutes waited} --- every value in between can happen};
\end{tikzpicture}
\end{center}
\\
& \notesprompt{Discrete or continuous? Say how you know.}
\begin{probgrid}{|Y|Y|}
\pcell{1}{Minutes a student took to finish a quiz}{2.0cm}{Continuous --- a better stopwatch splits 12 min into 12.4} &
\pcell{2}{Questions a student attempted on the quiz}{2.0cm}{Discrete --- 7 or 8 questions, never 7.5} \\ \hline
\end{probgrid}
\\
& \begin{probgrid}*{|Y|Y|}
\pcell{3}{Rainfall, read from a gauge that reports whole millimetres}{2.0cm}{Continuous --- the rain was a real amount; the gauge rounded it} &
\pcell{4}{Cars in the parking lot at noon}{2.0cm}{Discrete --- you count cars} \\ \hline
\end{probgrid}
\\ \hline
% ── ROW 2: displays that keep every value ────────────────────────────────────
\mainidea[Displays that keep]{Every value} &
Twenty-four students took a 50-minute quiz. How long each one took, shown two ways.
\textbf{Both pictures keep every individual value.}

\vspace{4pt}
\begin{minipage}[t]{0.38\linewidth}
\vspace{0pt}
\centering
\renewcommand{\arraystretch}{1.3}
{\normalsize
\begin{tabular}{@{} r | l @{}}
\textbf{1} & 1\,2\,2\,4\,5\,5\,6\,7\,8\,8\,9 \\
\textbf{2} & 0 \\
\textbf{3} & 4\,5\,6\,6\,8\,9 \\
\textbf{4} & 0\,1\,2\,3\,4\,6 \\
\end{tabular}}\\[3pt]
{\scriptsize Key: 3\,|\,4 means 34 minutes}\\[6pt]
This is a \labelbox{2.4cm}{stem plot}
\end{minipage}
\hfill
\begin{minipage}[t]{0.58\linewidth}
\vspace{0pt}
\centering
\begin{tikzpicture}[scale=0.82]
  \draw[->, thick] (-0.25,0) -- (8.4,0);
  \foreach \x/\lab in {0/10, 2/20, 4/30, 6/40, 8/50} {
    \draw (\x,-0.07) -- (\x,0.07); \node[below, font=\tiny] at (\x,-0.1) {\lab};}
  \foreach \v/\n in {11/1, 12/2, 14/1, 15/2, 16/1, 17/1, 18/2, 19/1, 20/1,
                     34/1, 35/1, 36/2, 38/1, 39/1, 40/1, 41/1, 42/1, 43/1, 44/1, 46/1} {
    \foreach \k in {1,...,\n} {
      \fill[navy] ({(\v-10)*0.2},{0.19*\k}) circle (0.065);}}
  \node[font=\tiny] at (4.0,-0.52) {Minutes to finish};
\end{tikzpicture}\\[6pt]
This is a \labelbox{2.4cm}{dotplot}
\end{minipage}
\\
& \notesprompt{Read them.}
\begin{probgrid}{|Y|Y|}
\pcell{5}{Fastest time and slowest time}{1.8cm}{11 minutes and 46 minutes} &
\pcell{6}{How many took exactly 36 minutes? Which picture makes that easiest to see?}{1.8cm}{2; the stem plot (the two 6s on the 3 row)} \\ \hline
\end{probgrid}
\\
& \begin{probgrid}*{|Y|Y|}
\pcell{7}{How many finished in under 20 minutes?}{1.8cm}{11 students} &
\pcell{8}{\textbf{In your own words:} what is the same about these two pictures?}{1.8cm}{Same 24 values, every one still visible in each} \\ \hline
\end{probgrid}
\\ \hline
% ── ROW 3: histograms ────────────────────────────────────────────────────────
\mainidea{Histograms} &
Twenty-four dots is readable; two hundred would not be. So we \textbf{group}:

\vspace{3pt}
\stepnum{1}\ Choose an \textbf{interval width} and mark the intervals on a number line.\par
\stepnum{2}\ Count the values that fall in each interval.\par
\stepnum{3}\ Draw one bar per interval, its height the count. \textbf{The bars touch.}

\vspace{5pt}
\textbf{9.} Count the 24 times into five-minute intervals:

\vspace{2pt}
\footnotesize\setlength{\tabcolsep}{2pt}%
\renewcommand{\arraystretch}{1.4}
\begin{tabularx}{\linewidth}{@{} l Y Y Y Y Y Y Y Y @{}}
\rowcolor{sky}
\textbf{Minutes} & 10--14 & 15--19 & 20--24 & 25--29 & 30--34 & 35--39 & 40--44 & 45--49 \\
\hline
\textbf{Students} & \ans{4} & \ans{7} & \ans{1} & \ans{0} &
\ans{1} & \ans{5} & \ans{5} & \ans{1} \\
\end{tabularx}\small
\\
& \begin{center}
\begin{tikzpicture}[scale=0.9]
  \foreach \y/\lab in {0/0, 1.4/5, 2.8/10} {
    \draw[linegray, very thin] (0,\y) -- (5.6,\y);
    \draw (-0.08,\y) -- (0.08,\y) node[left, font=\tiny, xshift=-2pt] {\lab};}
  \foreach \i/\c in {0/4, 1/7, 2/1, 3/0, 4/1, 5/5, 6/5, 7/1} {
    \fill[navy] ({\i*0.7},0) rectangle ({(\i+1)*0.7},{\c*0.28});
    \draw[white, very thin] ({\i*0.7},0) rectangle ({(\i+1)*0.7},{\c*0.28});}
  \draw[->, thick] (0,0) -- (0,3.9);
  \draw[->, thick] (0,0) -- (6.0,0);
  \foreach \x/\lab in {0/10, 1.4/20, 2.8/30, 4.2/40, 5.6/50} {
    \draw (\x,-0.08) -- (\x,0.08); \node[below, font=\tiny] at (\x,-0.12) {\lab};}
  \node[font=\tiny] at (2.8,-0.62) {Minutes to finish};
  \node[rotate=90, font=\tiny] at (-0.72,1.6) {Students};
\end{tikzpicture}
\end{center}
\\
& \textbf{That table is already the histogram.} What you give up is the individual values.
A \textbf{cumulative graph} counts how many fall \emph{at or below} each value instead; it returns
in Lesson~1.6.
\notesprompt{Use the histogram.}
\begin{probgrid}{|Y|Y|}
\pcell{10}{What does the height of the 35--39 bar tell you?}{2.2cm}{5 students finished in 35--39 minutes} &
\pcell{11}{Two students took exactly 36 minutes. Can this picture tell you that? What did grouping
cost?}{2.2cm}{No --- grouping hid the individual values; you can no longer see any single time} \\ \hline
\end{probgrid}
\\ \hline
% ── ROW 4: interval width — THE CRUX ─────────────────────────────────────────
\mainidea[The same 24 times,]{Drawn twice} &
Both pictures were built from the \emph{same} 24 times. Only the \textbf{interval width}
changed.

\vspace{2pt}
\begin{center}
\begin{tikzpicture}[scale=0.72]
  \foreach \y/\lab in {0/0, 1.4/5, 2.8/10} {
    \draw[linegray, very thin] (0,\y) -- (5.6,\y);
    \draw (-0.08,\y) -- (0.08,\y) node[left, font=\tiny, xshift=-2pt] {\lab};}
  \foreach \i/\c in {0/4, 1/7, 2/1, 3/0, 4/1, 5/5, 6/5, 7/1} {
    \fill[navy] ({\i*0.7},0) rectangle ({(\i+1)*0.7},{\c*0.28});
    \draw[white, very thin] ({\i*0.7},0) rectangle ({(\i+1)*0.7},{\c*0.28});}
  \draw[->, thick] (0,0) -- (0,4.0);
  \draw[->, thick] (0,0) -- (6.0,0);
  \foreach \x/\lab in {0/10, 1.4/20, 2.8/30, 4.2/40, 5.6/50} {
    \draw (\x,-0.08) -- (\x,0.08); \node[below, font=\tiny] at (\x,-0.12) {\lab};}
  \node[font=\tiny] at (2.8,-0.78) {Minutes to finish};
  \node[font=\bfseries\scriptsize] at (2.8,4.3) {Blocks of 5 minutes};
  \begin{scope}[xshift=8.0cm]
    \foreach \y/\lab in {0/0, 1.4/5, 2.8/10} {
      \draw[linegray, very thin] (0,\y) -- (5.6,\y);
      \draw (-0.08,\y) -- (0.08,\y) node[left, font=\tiny, xshift=-2pt] {\lab};}
    \foreach \i/\c in {0/12, 1/12} {
      \fill[goldacc] ({\i*2.8},0) rectangle ({(\i+1)*2.8},{\c*0.28});
      \draw[white, very thin] ({\i*2.8},0) rectangle ({(\i+1)*2.8},{\c*0.28});}
    \draw[->, thick] (0,0) -- (0,4.0);
    \draw[->, thick] (0,0) -- (6.0,0);
    \foreach \x/\lab in {0/10, 2.8/30, 5.6/50} {
      \draw (\x,-0.08) -- (\x,0.08); \node[below, font=\tiny] at (\x,-0.12) {\lab};}
    \node[font=\tiny] at (2.8,-0.78) {Minutes to finish};
    \node[font=\bfseries\scriptsize] at (2.8,4.3) {Blocks of 20 minutes};
  \end{scope}
\end{tikzpicture}
\end{center}
\\
& \fcolorbox{navy}{sky}{\parbox{\dimexpr\linewidth-2\fboxsep-2\fboxrule\relax}{\centering
\textbf{A histogram is a choice, not a fact.} Changing the interval width changes the shape
--- without changing a single data value.}}

\vspace{3pt}
\textbf{In your own words:} what changed between the two pictures, and what did not?
\writespace{1.6cm}{The width of the blocks changed. Not one of the 24 times changed.}
\\
& \notesprompt{Compare the two pictures.}
\begin{probgrid}{|Y|Y|}
\pcell{12}{Add up each picture's bars. Did either one drop anybody?}{2.4cm}{$4+7+1+0+1+5+5+1=24$ and $12+12=24$. Nobody was dropped.} &
\pcell{13}{Which picture shows the class splitting into a fast group and a slow group? What does
the other one say instead?}{2.4cm}{Blocks of 5. The other says students finished evenly across the period.} \\ \hline
\end{probgrid}
\\
& \begin{probgrid}*{|Y|Y|}
\pcell{14}{The teacher wants to know whether the quiz was too long. Which picture should she
use, and why?}{2.4cm}{Blocks of 5: the gap at 25--30 is what she needs to see; 20-minute blocks average it away} &
\pcell{15}{Redraw the times in blocks of \emph{one} minute. What would that picture look like,
and why is it not better either?}{2.4cm}{Almost every bar 0 or 1 tall --- every value back, no shape at all} \\ \hline
\end{probgrid}
\\ \hline
% ── ROW 5: guided practice — WE DO, together, fresh context ──────────────────
\mainidea[Guided practice]{The Bus Stop} &
Waiting times, in minutes, for 18 riders at one bus stop on a weekday morning.

\vspace{2pt}
\begin{center}
\begin{tikzpicture}[scale=1.0]
  \draw[->, thick] (-0.2,0) -- (7.4,0);
  \foreach \x/\lab in {0/0, 1/2, 2/4, 3/6, 4/8, 5/10, 6/12, 7/14} {
    \draw (\x,-0.07) -- (\x,0.07); \node[below, font=\tiny] at (\x,-0.1) {\lab};}
  \foreach \x/\n in {0.5/2, 1/3, 1.5/4, 2/3, 2.5/2, 3/1, 5.5/1, 6/2} {
    \foreach \k in {1,...,\n} { \fill[navy] (\x,{0.22*\k}) circle (0.08); }}
  \node[font=\tiny] at (3.6,-0.52) {Minutes waited};
\end{tikzpicture}
\end{center}
\\
& \notesprompt{We work these together.}
\begin{probgrid}{|Y|Y|}
\pcell{16}{How many riders are shown, and how many waited longer than 8 minutes?}{2.2cm}{18 riders; 3 waited longer than 8 minutes} &
\pcell{17}{Is waiting time discrete or continuous? Say how you know from the variable, not the
picture.}{2.2cm}{Continuous --- time is measured; a finer clock splits 3 min} \\ \hline
\end{probgrid}
\\
& \begin{probgrid}*{|Y|Y|}
\pcell{18}{Count the times into 5-minute blocks (0--4, 5--9, 10--14). In that picture, do the
two riders at 11 and 12 minutes still stand apart? Why?}{2.6cm}{0--4: 12, 5--9: 3, 10--14: 3. No --- the last two bars are equal, so the empty stretch at 7--10 vanished} &
\pcell{19}{Which picture would you show the bus company, and what would you tell them?}{2.6cm}{The dotplot: most wait under 6 minutes, but two waited 11--12 --- that gap is the story} \\ \hline
\end{probgrid}
\\ \hline
\end{guidednotes}

% ── THE NOTES END AT GUIDED PRACTICE ─────────────────────────────────────────
% There is deliberately no "you do" here: ../homework/ is the individual practice,
% started in class in the last 8 minutes while the teacher circulates.

\end{document}
